% Generated draft: requires mathematical and source review.
\documentclass[11pt]{article}
\usepackage[margin=25mm]{geometry}
\usepackage{fontspec}
\setmainfont{Latin Modern Roman}
\usepackage{amsmath,amssymb,mathtools,unicode-math}
\setmathfont{Latin Modern Math}
\usepackage{enumitem,needspace,xurl,hyperref}
\setlength{\parindent}{0pt}
\setlength{\parskip}{5pt}
\setlength{\emergencystretch}{3em}
\newcommand{\N}{\mathbb{N}}
\newcommand{\Z}{\mathbb{Z}}
\newcommand{\Q}{\mathbb{Q}}
\newcommand{\R}{\mathbb{R}}
\newcommand{\C}{\mathbb{C}}
\newcommand{\F}{\mathbb{F}}
\newcommand{\A}{\mathbb{A}}
\newcommand{\PP}{\mathbb{P}}
\newcommand{\E}{\mathbb{E}}
\newcommand{\eps}{\varepsilon}
\newcommand{\dd}{\,\mathrm{d}}
\begin{document}
\section*{UnsolvedMath ID: 800985}
\textbf{Source problem code:} catalogue entry
\subsection{Definitions and mathematical statement}
In the usual differential-geometric interpretation, a projective structure \(\mathcal P\) on a manifold \(M\) is an equivalence class of torsion-free affine connections having the same unparameterized geodesics. Equivalently, representatives \(\nabla\) and \(\widehat\nabla\) are projectively equivalent if there is a one-form \(\alpha\) such that
\[
\widehat\nabla_XY=\nabla_XY+\alpha(X)Y+\alpha(Y)X
\]
for all vector fields \(X,Y\). A metric \(g\) on an open subset \(U\subseteq M\) is compatible with \(\mathcal P\) if its Levi--Civita connection belongs to \(\mathcal P|_U\). Thus, for a representative \(\nabla\) of \(\mathcal P|_U\), compatibility means that some one-form \(\alpha\) on \(U\) satisfies
\[
\nabla^g_XY=\nabla_XY+\alpha(X)Y+\alpha(Y)X
\qquad\text{for all vector fields }X,Y\text{ on }U.
\]
These definitions explain the conventional connection-based meaning of the terminology; the supplied record does not explicitly specify this convention, the dimension or regularity of \(M\), or the permitted metric signatures.

\textbf{Requested construction (Burns--Matveev, Problem 66).} Construct a system of scalar invariants \(\{I_\lambda\}_{\lambda\in\Lambda}\) of a projective structure whose vanishing is necessary and sufficient for local compatibility with a metric in a neighborhood of almost every point. A scalar invariant is understood here as a scalar quantity intrinsically determined by the projective structure, independent of coordinates and of the representative connection. The requested equivalence has the schematic form
\[
\bigl[\text{all }I_\lambda(\mathcal P)\text{ vanish}\bigr]
\quad\Longleftrightarrow\quad
\left[
\begin{gathered}
\text{for almost every }p\in M,\text{ there exist an open neighborhood }U\ni p\\
\text{and a metric }g\text{ on }U\text{ compatible with }\mathcal P|_U.
\end{gathered}
\right]
\]
The neighborhoods and metrics may depend on \(p\); no global compatible metric is requested. This is a qualified formulation: the supplied record does not define ``almost every,'' specify where or in what sense the invariants must vanish, or state whether the system must be finite or subject to a differential-order bound. It also leaves open whether ``metric'' requires positive definiteness or allows other nondegenerate signatures. These choices must be resolved before the schematic equivalence becomes a fully specified mathematical statement.
\subsection{Short English statement}
Construct scalar invariants that detect whether a projective structure locally has the same unparameterized geodesics as a metric near almost every point. The metric class, ambient dimension and regularity, meaning of ``almost every,'' and precise vanishing condition remain unspecified in the supplied source record.
\subsection{Sources}
\begin{itemize}
\item Local UnsolvedMath registry, permanent ID \(800985\), source code \texttt{catalogue entry}, title ``Open Problems and Questions about Geodesics --- Problem 66.'' The supplied record attributes the problem to Burns and Matveev, with proposed year 2019.
\item Burns and Matveev, \emph{Open Problems and Questions about Geodesics}, source list dated 2019 in the supplied provenance, Problem 66. Supplied source URL: \url{https://arxiv.org/abs/1308.5417}. The record reports access on 2026-10-16 and extraction by \texttt{arxiv\_latex}; these are supplied provenance claims, not independently verified access or extraction.
\item The local record labels the problem \texttt{open}. This is a supplied registry status, not an independently verified literature assessment. Supplied upstream revision: \texttt{d371ee917ac9647fbb06cd6fc529deec6f3dda82}.
\end{itemize}
\end{document}
